% =====================================================================
\chapter{Puesta en producción en planta}
\label{tema:6}
% =====================================================================

Un controlador que funciona en el gemelo digital está a mitad de camino. En planta hay operadores que deben entender y poder anular lo que hace, alarmas que no deben saturarlos, un sistema de control básico y otro de seguridad que no se pueden comprometer, una red industrial que es un objetivo de ciberataques y un marco normativo. Este tema recorre esos requisitos. La arquitectura por capas sigue el capítulo~15 («Architecture and System Design») de Åström y Murray~\cite{am}; la gestión de alarmas, la norma ISA-18.2~\cite{isa182}; la seguridad funcional, IEC~61511~\cite{iec61511}; la ciberseguridad, la guía NIST SP~800-82~\cite{nist} y la serie IEC~62443~\cite{iec62443}; la validación y el seguimiento de modelos, el curso CS329S de Stanford~\cite{cs329s} y Sculley \emph{et al.}~\cite{sculley}, y la regulación, el Reglamento europeo de inteligencia artificial~\cite{aiact}. Las normas son documentos de pago: aquí se usa solo lo que sus organismos difunden públicamente.

\section{Arquitectura por capas}

\addtocounter{figure}{1}%

La \cref{fig:capas} resume el principio de diseño: cada capa trabaja en una escala de tiempo y confía en la inferior. El MPC o la política de RL \emph{no} mueven válvulas: envían consignas a los PID existentes, que siguen funcionando si la capa superior falla. Entre ambas, un filtro valida cada consigna. El \emph{sistema instrumentado de seguridad} (SIS), que detiene el proceso ante condiciones peligrosas, es independiente de todo lo anterior, con sus propios sensores y lógica, y se diseña según IEC~61511~\cite{iec61511}, la aplicación al sector de procesos de la seguridad funcional de IEC~61508, que asigna a cada función de seguridad un nivel de integridad (SIL, de 1 a 4). Ningún componente de aprendizaje automático forma parte del SIS. El caso de Google~\cite{gamblegao} sigue el mismo esquema: la IA en la nube propone, el sistema local verifica y aplica, y las reglas clásicas quedan como respaldo.

\section{Modos de operación y lazo humano}

\begin{definicion}[Modos de operación]\label{def:modos}
Se distinguen tres modos de creciente autonomía: \emph{asesor} (el sistema recomienda y el operador decide y ejecuta), \emph{supervisado} (el sistema ejecuta dentro de límites y el operador vigila y puede intervenir en cualquier momento) y \emph{autónomo con supervisión} (el sistema opera con salvaguardas automáticas y el operador interviene por excepción). En todos ellos el operador puede volver al control básico.
\end{definicion}

Un despliegue prudente recorre los modos en orden y solo avanza cuando hay evidencia: el caso de los centros de datos pasó de recomendaciones en 2016~\cite{evansgao} a control autónomo con supervisión en 2018~\cite{gamblegao}. El paso entre modos no debe provocar saltos en los actuadores.

\begin{proposicion}[Transferencia sin saltos]\label{prop:bumpless}
Sea un PI en forma posicional, $u_k=K_ce_k+I_k$, que toma el control en el instante $k_0$ de un sistema que aplicaba $u_{k_0-1}$. Si se inicializa $I_{k_0}=u_{k_0-1}-K_ce_{k_0}$, entonces $u_{k_0}=u_{k_0-1}$.
\end{proposicion}

Con la zona del \cref{ej:caso-pi}, si la capa supervisora aplicaba $1{,}2$ kW y el error al recuperar el PI es de $0{,}4$ K, el integrador se inicializa en $1{,}2-1{,}429\cdot0{,}4=0{,}629$ kW. El mismo principio se aplica al pasar del PI al MPC: el MPC debe penalizar los cambios de la acción respecto a la última aplicada.

\section{Alarmas}

La norma ANSI/ISA-18.2, base de la internacional IEC~62682~\cite{isa182}, organiza la gestión de alarmas en un ciclo de vida (filosofía, identificación, racionalización, diseño, implantación, operación, mantenimiento, seguimiento, gestión del cambio y auditoría) y da referencias de carga: alrededor de 150 alarmas al día (una cada 10 minutos de media) se considera muy probablemente aceptable y alrededor de 300 al día (dos cada 10 minutos), el máximo manejable. Una alarma solo debe existir si exige una acción del operador. Una capa de IA añade fuentes nuevas que hay que racionalizar con el mismo criterio: residuos del modelo fuera de banda, proximidad a restricciones, salidas a modo de respaldo, pérdida de comunicación con la capa supervisora o predicciones con incertidumbre alta.

\section{Validación y despliegue}

\begin{definicion}[Etapas de validación]\label{def:validar}
(i) \emph{Validación fuera de línea}: el controlador se evalúa en el gemelo y sobre históricos, con escenarios no usados en su diseño. (ii) \emph{Modo sombra}: el controlador calcula sus decisiones con datos en tiempo real, pero no las aplica; se comparan con las del sistema existente y se comprueba que serían seguras. (iii) \emph{Experimento en planta}: se alterna el controlador nuevo y el existente por periodos (días o semanas) y se mide la diferencia con un diseño de comparación pareada~\cite[cap.~2]{montgomery}. (iv) \emph{Despliegue con seguimiento}: se vigilan los residuos del modelo, la distribución de las entradas y los indicadores, con criterios de vuelta atrás definidos de antemano.
\end{definicion}

\addtocounter{definicion}{1}%

Los modelos degradan: cambian los equipos, el uso del edificio o el producto, y las relaciones aprendidas dejan de valer. Sculley \emph{et al.}~\cite{sculley} describen cómo un sistema de aprendizaje automático acumula «deuda técnica» por dependencias de datos ocultas, bucles de realimentación (el controlador cambia los datos con los que se reentrenará) y configuración, y las clases~10 y~11 de CS329S~\cite{cs329s} tratan la detección de cambios de distribución y el seguimiento. En control, la señal de seguimiento más útil es el residuo del modelo en simulación libre sobre ventanas recientes: un sesgo sostenido indica que hay que recalibrar el gemelo.

\section{El filtro de seguridad}

\begin{definicion}[Filtro de seguridad]\label{def:filtro}
Sea un modelo $x_{k+1}=f(x_k,u_k,w_k)$ con perturbación acotada $w_k\in W$ y un conjunto seguro $\Sset_{\mathrm{seg}}$. Para cada $x$ se define el conjunto de acciones seguras $\Uset_{\mathrm{seg}}(x)=\{u\in\Uset:\ f(x,u,w)\in\Sset_{\mathrm{seg}}\ \forall w\in W\}$. El \emph{filtro de seguridad} sustituye la acción propuesta $u^{\mathrm{p}}$ por su proyección $u=\proy_{\Uset_{\mathrm{seg}}(x)}(u^{\mathrm{p}})$.
\end{definicion}

\begin{proposicion}[Invariancia bajo el filtro]\label{prop:invariancia}
Si $\Uset_{\mathrm{seg}}(x)\ne\emptyset$ para todo $x\in\Sset_{\mathrm{seg}}$ y $x_0\in\Sset_{\mathrm{seg}}$, entonces, cualquiera que sea la política que propone las acciones, $x_k\in\Sset_{\mathrm{seg}}$ para todo $k$ y toda secuencia de perturbaciones en $W$.
\end{proposicion}

La garantía es tan buena como el modelo y la cota de las perturbaciones: el filtro usa un modelo sencillo, conservador y verificable, no la red del agente. Cuando $\Uset_{\mathrm{seg}}$ es un intervalo, la proyección es un simple recorte. Se completa con límites de velocidad de cambio (Lazic \emph{et al.}~\cite{lazic} limitaban la variación de cada actuador entre pasos), un \emph{perro guardián} (\emph{watchdog}) que devuelve el control a los PID si la capa supervisora deja de enviar consignas válidas, y la transferencia sin saltos de la \cref{prop:bumpless}.

\addtocounter{definicion}{1}%

\section{Ciberseguridad de la tecnología de operación}

Los sistemas de control industrial (tecnología de operación, OT) priorizan la disponibilidad y la seguridad física por encima de la confidencialidad, funcionan durante décadas con equipos difíciles de actualizar y usan protocolos diseñados sin autenticación. La guía NIST SP~800-82, revisión~3, de 2023~\cite{nist}, adapta la gestión de riesgos de ciberseguridad a este entorno. La serie IEC~62443~\cite{iec62443}, desarrollada por IEC e ISA, se organiza en cuatro bloques (general; políticas y procedimientos; sistema; componentes) y aporta dos conceptos centrales: la partición del sistema en \emph{zonas} de activos con requisitos de seguridad comunes, comunicadas solo por \emph{conductos} definidos, y los \emph{niveles de seguridad} SL~1 a SL~4, ligados a la capacidad del adversario que se quiere resistir, desglosados en siete requisitos fundamentales (identificación y autenticación, control de uso, integridad del sistema, confidencialidad de los datos, flujo de datos restringido, respuesta oportuna a eventos y disponibilidad de recursos).

Para una capa de IA, de ahí se derivan reglas prácticas: (i) la capa supervisora vive en su propia zona y solo escribe consignas por un conducto autenticado y registrado, nunca acceso directo a los controladores; (ii) el sistema local valida cada consigna (el filtro de seguridad es también una defensa: limita el daño de una consigna maliciosa); (iii) si el modelo se ejecuta en la nube, la planta debe poder operar sin ella; (iv) los datos de entrenamiento son un activo: un atacante que los manipule puede inducir un modelo que se comporte mal en condiciones concretas, y (v) actualizar un modelo es un cambio en el sistema de control y pasa por la gestión del cambio.

\section{Normas y regulación}

\begin{table}[htbp]
\centering\small
\begin{tabularx}{\textwidth}{@{}l>{\raggedright\arraybackslash}X>{\raggedright\arraybackslash}X@{}}
\toprule
Referencia & Ámbito & Qué implica para una capa de IA \\
\midrule
IEC 61508 / IEC 61511~\cite{iec61511} & seguridad funcional; sistemas instrumentados de seguridad en procesos & la IA queda fuera del SIS; el análisis de riesgos debe considerar sus modos de fallo \\
ANSI/ISA-18.2, IEC 62682~\cite{isa182} & gestión de alarmas & racionalizar las alarmas nuevas; vigilar la carga del operador \\
IEC 62443~\cite{iec62443}, NIST SP 800-82~\cite{nist} & ciberseguridad de sistemas de automatización & zonas y conductos; autenticación y registro de consignas \\
Reglamento (UE) 2024/1689~\cite{aiact} & inteligencia artificial & los sistemas destinados a ser componentes de seguridad en la gestión y operación de infraestructuras críticas, como el suministro de agua, gas, calefacción o electricidad, se clasifican como de alto riesgo (anexo~III, punto~2) \\
\bottomrule
\end{tabularx}
\caption{Marco normativo de referencia.}
\label{tab:normas}
\end{table}

Los sistemas de alto riesgo del Reglamento de IA deben cumplir los requisitos de la sección~2 de su capítulo~III (artículos~8 a~15), entre ellos la gestión de riesgos, la gobernanza de los datos, el registro de eventos, la supervisión humana (artículo~14: el sistema debe diseñarse para que las personas puedan entender sus límites, detectar anomalías, no confiar en exceso en sus salidas, anular sus decisiones y detenerlo de forma segura) y la precisión, robustez y ciberseguridad. Un optimizador de climatización de oficinas no suele ser un componente de seguridad; el control de una red de calefacción urbana o de una subestación puede serlo. La clasificación de cada caso es una cuestión jurídica que debe evaluarse con especialistas; lo que el módulo aporta es que la arquitectura de la \cref{fig:capas} facilita cumplir esos requisitos.

\begin{table}[htbp]
\centering\footnotesize
\begin{tabularx}{\textwidth}{@{}>{\raggedright\arraybackslash}p{3.1cm}>{\raggedright\arraybackslash}X>{\raggedright\arraybackslash}p{2.0cm}>{\raggedright\arraybackslash}X@{}}
\toprule
Herramienta & Para qué sirve & Licencia o tipo & Cuándo elegirla \\
\midrule
OPC UA & comunicación industrial con seguridad (autenticación, cifrado) & estándar abierto & leer medidas y escribir consignas desde la capa supervisora \\
Historiadores de proceso & almacenamiento de series temporales de planta & comercial o libre & datos de identificación y de seguimiento \\
MLflow u otros registros de modelos & versionado de modelos, datos y métricas & libre (Apache-2.0) & trazabilidad de qué modelo controló la planta y cuándo \\
Contenedores y orquestación & despliegue reproducible de la capa supervisora & libre & separar entornos de prueba y producción \\
Evidently, WhyLogs & seguimiento de distribuciones y de la deriva de los datos & libres & alarmas de degradación del modelo \\
\bottomrule
\end{tabularx}
\caption{Tecnologías para la puesta en producción.}
\label{tab:tec6}
\end{table}

\addtocounter{definicion}{1}%

